\section{Introducción}
El fútbol autónomo combina estados físicos continuos, interacción entre jugadores
y recompensas dispersas. Resolver habilidades por separado no asegura un ataque
completo eficaz: el estado que deja una habilidad condiciona el desempeño de la
siguiente. Se aborda este problema en RoboCup 2D mediante un coordinador que
selecciona persecución, conducción o definición.

El estudio compara aprendizaje de valor con DQN/Double DQN y métodos actor-crítico
con A2C/PPO, tomando el control tabular como referencia. Se integra PBRS para
orientar la exploración y BC/DAgger para inicializar las políticas antes del
ajuste con RL. La evaluación distingue aprendizaje por habilidad, efecto del
shaping y desempeño de la jerarquía frente a una política plana. El objetivo es
alcanzar al menos 80\,\% de goles en 200 episodios de prueba, examinando también
la eficiencia muestral y las limitaciones de generalización. Las redes y los
bucles de aprendizaje se implementan en PyTorch.
